\chapter{Introducción}
\label{chap:introduccion}

Las aplicaciones Ionic desarrolladas con Capacitor suelen probarse en Node.js con frameworks como Jest o Jasmine. En ese entorno los imports \texttt{@capacitor/*} no reproducen el comportamiento nativo del dispositivo: si no se sustituyen por implementaciones alineadas al contrato del plugin, la suite falla o ejercita sólo caminos superficiales. La \textbf{cobertura de líneas y ramas} pierde entonces utilidad, porque gran parte del código que depende de APIs nativas no llega a ejecutarse de manera significativa en el runner. La documentación oficial orienta a mockear plugin por plugin; trabajos recientes sobre cobertura y generación en JavaScript \cite{schafer2024testpilot,olsthoorn2024syntest_javascript} dan marcos generales, pero no resuelven por sí solos la combinación bridge Capacitor + entorno Node reproducible.

La propuesta integra un \textbf{módulo de simulación de plugins nativos Capacitor} (respuestas configurables y \textit{test doubles} centralizados, acoplables a Jest y/o Jasmine), un \textbf{procedimiento para mejorar cobertura de líneas y ramas} mediante escenarios explícitos de éxito, error y permisos---sin equiparar esa medición con pruebas en hardware---y, en apoyo operativo, una \textbf{capa MCP auxiliar} que automatice lo repetible del ciclo de medición (suite con cobertura, lectura del reporte, Sonar u otros chequeos estáticos en solo lectura cuando existan). La evaluación será \textbf{cuantitativa y comparativa}: mismos comandos y versiones documentadas entre una \textbf{línea base} previa y el estado con el módulo integrado, registrando además el tiempo dedicado a preparar mocks frente al uso del simulador.

Los capítulos siguientes formulan el problema, la justificación, los objetivos, el alcance y el diseño experimental. Toda afirmación sobre cobertura debe explicitar que corresponde a unitarias en Node bajo simulación, no a pruebas de integración en dispositivo.
